\section{Weekly Summary}

\begin{tcolorbox}[colback=yellow!10,colframe=black!50,title=AI-Generated Content]
\noindent The summary below was written by Claude (Anthropic) from that week's regression outputs and series descriptions. It has not been reviewed or fact-checked by a human, and should be read as an automated, informal recap -- not analysis.
\end{tcolorbox}

Welcome back to another thrilling installment of Pulling Two Random Numbers Out of a Government Database and Shoving Them Into a Line. As always, the premise is as noble as it is pointless: every day we yank a fresh pair of unrelated economic series from FRED, draw a straight line between them, and then act surprised when the line fits. The whole project is a monument to spurious correlation, so please do not use anything below to, say, run a hedge fund or win an argument with your uncle.

We kicked off Sunday with goods-producing employment in Springfield, Massachusetts versus a nursing home price index, which posted a gaudy raw R-squared of 0.78 -- the kind of number that makes you want to pop champagne. Then we detrended it and the relationship evaporated into literally nothing (R-squared 0.02, p-value basically 0.52, which is statistics-speak for "lol no"). Textbook trend mirage: both series drifted over the years, the regression noticed the drift, and we almost fell for it. Monday's IT manufacturing orders versus Missouri bank equity returns was more honest about being weak -- a modest raw fit that limped down but at least stayed technically significant after detrending, which is the equivalent of a participation trophy. Tuesday pitted East North Central residential electricity prices against restaurant sales, and here's where things got mildly interesting: the raw fit was strong (0.74), and even after stripping the trends out it held on to a respectable 0.21 with a p-value that refuses to be ignored. I'm not saying your dinner bill is powered by the electric company, but I'm also not not saying it.

The real headline-grabbers came midweek. Wednesday's matchup -- Puerto Rico covered employment versus U.S. travel service exports -- was the inverse of the usual gag: the raw relationship was a shrug (R-squared 0.06, not even significant), but detrending the thing somehow cranked it UP to 0.52 with a p-value sporting eight zeros. A relationship that only appears once you remove the trends is a genuinely rare and spooky little creature, and I have no earthly idea what it means. Friday gave us the week's crown jewel: Norwegian recession indicators versus Northeast adjustable-rate mortgage margins, which held strong in both the raw AND detrended versions (R-squared climbing from 0.17 to 0.22, p-values so small they need scientific notation to even exist). Two discontinued series from different countries, both dead for years, reaching out from the grave to hold hands. Beautiful. Meaningless, almost certainly, but beautiful.

Rounding out the week, Thursday matched the trade balance against Midwest chicken breast prices, posting a solid raw fit that survived detrending in weakened form -- a mild "huh" rather than a full gasp. And Saturday closed us out with West Coast banana prices versus Mideast machinery manufacturing output, which looked promising raw (0.45) and then face-planted completely once detrended (0.02, p-value 0.52), joining the Springfield special in this week's Hall of Trend Mirages. So the final tally: a couple of fun survivors, a couple of spectacular collapses, and absolutely zero reason to believe any of it. See you next week for more rigorous nonsense.
